\section{Introduction}
\label{sec:intro}

\IEEEPARstart{C}{ontinuous} and reliable Earth Observation (EO) using spaceborne optical remote sensing systems is fundamental to modern environmental monitoring, agricultural yield estimation, natural disaster assessment, and urban land-use governance. High-resolution optical sensors, such as the Indian Space Research Organisation's (ISRO) Resourcesat-2/2A Linear Imaging Self-Scanning Sensor (LISS-IV), capture multispectral imagery at an exquisite spatial ground sampling distance (GSD) of $5.0\,\text{m}$. Such data provides unprecedented spectral fidelity in visible and near-infrared (VNIR) bands, enabling fine-scale vegetation canopy modeling and parcel-level crop classification.

Despite these advanced orbital sensing capabilities, the utility of optical satellite imagery is severely compromised by persistent atmospheric cloud occlusions and accompanying cloud-shadow attenuation. Climatological studies indicate that approximately $66\%$ of the Earth's terrestrial surface is obscured by cloud cover at any given observation window, with cloud cover exceeding $85\%$ across tropical and subtropical regions during prolonged summer monsoon cycles (e.g., the South Asian monsoon from June to September). Dense cumulus clouds entirely extinguish optical reflectance via Mie scattering, while semi-transparent cirrus clouds introduce non-linear spectral distortion and diffuse haze. Consequently, operational downstream analytics suffer from extensive temporal data gaps, rendering real-time agricultural monitoring and flood extent mapping unreliable when most critically needed.

\begin{figure*}[!t]
\centering
\includegraphics[width=0.98\textwidth]{cloudfree_vision_v2_architecture.png}
\caption{\textbf{End-to-End Architectural Pipeline of RelieF-CR.} The framework integrates four heterogeneous multi-modal streams: (1) Cloudy LISS-IV optical imagery $[B, 3, 256, 256]$, (2) Sentinel-1 dual-polarization (VV/VH) SAR $[B, 2, 256, 256]$, (3) Sentinel-2 dry-season clear-sky temporal reference $[B, 3, 256, 256]$, and (4) CartoDEM 4-channel topographic surface descriptors $[B, 4, 256, 256]$. Dedicated convolutional stems extract intermediate feature representations at half-resolution ($128\times 128$). Learned Spatial Transformer Networks (STNs) dynamically compensate for sub-pixel inter-sensor registration discrepancies between SAR/temporal features and the optical reference. The aligned auxiliary features are concatenated and fused with optical query tokens via Windowed Multi-Head Cross-Attention over localized $8\times 8$ windows. A multi-scale decoder with residual skip connections generates the primary reconstructed clear-sky optical image $\hat{\mathbf{y}} \in [-1, 1]$ and an uncertainty log-variance map $\log \sigma^2 \in \mathbb{R}$. Training is driven by a comprehensive 7-term loss suite incorporating frequency-domain FFT and spectral consistency constraints.}
\label{fig:main_architecture}
\end{figure*}

To overcome optical cloud opacity, multi-sensor data fusion has emerged as the preeminent paradigm in Earth observation. Spaceborne Synthetic Aperture Radar (SAR) systems, such as the European Space Agency's (ESA) Copernicus Sentinel-1 constellation operating at C-band ($5.405\,\text{GHz}$), emit microwave radiation that completely penetrates cloud droplets, aerosol plumes, and rain cells irrespective of solar illumination or atmospheric conditions. Sentinel-1 provides dual-polarization backscatter measurements (VV and VH) that capture surface roughness, moisture content, and dielectric properties. However, SAR backscatter and optical multispectral reflectance represent fundamentally different physical interaction regimes. Radar backscatter is governed by surface geometric roughness, volumetric canopy scattering, and radar incidence angles, whereas optical reflectance measures electronic and vibrational absorption of solar photons across cellular chloroplasts and soil minerals. Furthermore, SAR imagery is intrinsically corrupted by coherent speckle noise, geometric layover, and foreshortening in complex terrain, making direct SAR-to-optical translation an ill-posed inverse problem.

To regularize this cross-modal synthesis, researchers have explored incorporating historical clear-sky optical passes as temporal priors (e.g., dry-season Sentinel-2 references). Nevertheless, temporal references alone fail during rapid land-cover transitions, such as seasonal crop growth, harvesting, or sudden flood inundation. Moreover, complex topographic relief significantly modulates both SAR backscatter (via local incidence angle variations) and optical illumination (via terrain self-shadowing), yet topographic priors remain almost entirely neglected in existing deep learning cloud removal architectures.

In this work, we propose \textbf{RelieF-CR} (Relief-Guided Cloud Removal), a unified, multi-modal cross-attention generative framework that synergistically fuses four distinct sensor modalities: high-resolution ($5.0\,\text{m}$) ISRO LISS-IV optical imagery, Sentinel-1 C-band dual-pol SAR, Sentinel-2 dry-season temporal references, and CartoDEM 4-channel topographic surface descriptors (elevation, slope, sin-aspect, cos-aspect). 

The primary contributions of this work are summarized as follows:
\begin{enumerate}
    \item \textbf{Heterogeneous Quad-Modal Fusion:} We formulate a cross-modal architecture that simultaneously conditions reconstruction on high-resolution optical, dual-pol SAR, temporal clear-sky references, and high-precision DEM surface descriptors, resolving the fundamental underdetermination of SAR-only inpainting.
    \item \textbf{Learned Spatial Alignment (STN):} We incorporate lightweight Spatial Transformer Networks (STNs) within the generator to dynamically estimate affine transformation fields, rectifying spatial registration jitter and resolution mismatches between SAR/temporal features and the optical coordinate grid prior to attention fusion.
    \item \textbf{Windowed Multi-Head Cross-Attention:} We design a localized windowed cross-attention mechanism where cloudy optical features serve as queries to selectively attend to auxiliary SAR, temporal, and topographic keys and values, dynamically extracting structural details exclusively over cloud-obscured regions.
    \item \textbf{Calibrated Heteroscedastic Uncertainty Quantification:} Unlike conventional deterministic models, our generator incorporates a dual-head output predicting both the reconstructed reflectance $\hat{\mathbf{y}}$ and a pixel-wise log-variance map $\log \sigma^2$, trained via Gaussian Negative Log-Likelihood (NLL) to provide reliable, calibrated confidence intervals.
    \item \textbf{Comprehensive 7-Term Loss Formulation:} We formulate a loss suite combining Gaussian NLL, multi-scale L1, spectral-normalized PatchGAN adversarial loss, feature matching, frequency-domain FFT magnitude loss, Sobel edge constraints, and NDVI spectral consistency, eliminating oversmoothing artifacts.
    \item \textbf{Rigorous Benchmark & Diagnostic Auditing:} Evaluated across 1,611 held-out test patches, our model achieves state-of-the-art performance (\textbf{29.10\,dB PSNR}, \textbf{0.963 SSIM}, \textbf{1.49$^\circ$ SAM}, \textbf{7.89 ERGAS}). We perform an exhaustive SAR-optical gradient correlation audit, mathematically proving that our model's structural correlation with SAR ($0.4791$) matches real optical ground truth ($0.4766$), ruling out radar over-imprinting artifacts.
\end{enumerate}
